\section{A path gap on every six-pole of girth at least seven}
\label{sec:root-six-planarity}

This observation applies at every rank.
It strengthens the bounded-shore reduction, but does not by itself prove a six-cut bound.

\begin{lemma}[Planarity of a saturated pole]
\label{lem:root-saturated-pole-planar}
Let $X$ be a connected normalized $k$-pole, with two ports $s,t$.
If $P_{st}(X)=2^{\beta(X)}$, then $X$ is planar.
\end{lemma}

\begin{proof}
Suppress all degree-two ports except $s,t$.
A loop cannot arise: it would be an off-route cycle, which makes the normalized terminal-path count at most one half.
The result is a saturated normalized two-pole of the same rank.
Apply Lemma~\ref{lem:full-two-pole-classification}.

Each bridgeless piece is a digon or a capped permutation ladder.
The permitted ladder blocks are singletons and adjacent transpositions.
Complete each block by adding an edge between its two left attachments, and another between its two right attachments.
For a singleton, the completed multigraph has two vertices and three parallel edges.
For a transposition, it is $K_4$.
Both are planar.

A planar two-edge join preserves planarity.
Indeed, put one designated edge of each planar graph on its outer face, remove these edges, and join corresponding endpoints across the two disks.
Apply this operation to successive completed blocks.
At the two ends, subdivide the remaining completion edges to form the terminal caps.
This constructs the capped ladder as a planar graph.
Joining planar pieces by bridges also preserves planarity.
Finally, restore the suppressed ports by subdivision.
\end{proof}

\begin{proposition}[Uniform path gap in a six-pole of large girth]
\label{prop:root-six-girth-path-gap}
Let $X$ be a connected normalized six-pole of girth at least seven and rank $h$.
Every terminal pair satisfies
\[
 P_{st}(X)\le\frac{15}{16}\,2^h.
\]
\end{proposition}

\begin{proof}
If some pair were saturated, Lemma~\ref{lem:root-saturated-pole-planar} would make $X$ planar.
Let $n,e,f$ denote its numbers of vertices, edges, and faces in a planar embedding.
The degree sum gives $2e=3n-6$.
The graph has minimum degree two and contains a cycle.
Every facial boundary walk therefore has length at least seven, even when bridges occur.
Thus $2e\ge7f=7(2-n+e)$, or $5e\le7n-14$.
Combining the two inequalities gives $n\le2$.
This contradicts the six distinct ports.
The terminal-path saturation gap now proves the assertion.
\end{proof}

The facial-walk assertion uses minimum degree two.
A facial walk with no cycle in its support would force a leaf in its finite tree support.
At such a leaf, the facial walk turns back immediately, so the original graph has degree one there.
This is impossible.
Equivalently, apply the usual planar girth inequality after deleting bridges and summing over the cyclic blocks.

\begin{remark}[Scope]
The argument does not claim that every planar six-pole is saturated.
It only excludes saturation when the girth is at least seven.
The preceding six-cycle-shore theorem already handles the smaller-girth cores relevant to the cycle-count problem.
\end{remark}
